\subsection{A retained-total complex motif at ground size twenty-four}
\label{sec:retained-complex-construction}
Use $h=24$, $v=\binom h3=2024$, $N=v^3$, and $m=h^3=13824$,
independently of the retained paired-bit arity $125000$.
The scalar domain is $\mathbb Z[i,1/2]$ and the binary label space is
$F=\mathbb F_2^h$ with its ordinary dot product. Triple indicators $t_T$
have norm one. The full space is $\mathcal F=F^{\otimes3}$.
All networks, labels and bases below are fixed finite data.

\paragraph{Side correction by shared source trees.}
The central scalar coefficient is $(|S\cap T|-1)/2$.
Its side correction is $+1/2$ when $S,T$ are disjoint and $-1/2$ when
$|S\cap T|=2$, and zero otherwise. Central plus side is the identity.
For the disjoint branch, recursively partition ordered disjoint triple pairs
into rectangles. On $D(a,b)$, use a row-star partition, a column-star partition,
or split the ground set and partition each source/target cardinality case by
the product of its two smaller disjointness partitions. Each case is disjoint,
and the selected recursive construction is an exact partition, not a cover.
For a rectangle with $a$ source triples and $b$ targets, aggregate its source
values by a balanced tree and retain one output use per target. It costs
$a+b-1$ roles. A singleton source retains its triple-line label.
For a larger source family use coordinate labels of its source unions.
If its full source union is the complete complement of any target, split that
rectangle into singleton-edge rectangles to avoid an alternating output
residual. Every repaired rectangle in the selected $h=24$ plan has one target,
so this repair preserves its role count. The repaired disjoint branch costs
$329811$ roles, with $123855$ additions and $205956$ outputs.
Separately retain the lexicographic source tree, grouped by least and second
point, only for its global total below.
For each fixed pair $P$, build a balanced tree on $x_{P\cup\{k\}}$, $k\notin P$.
For target $P\cup\{r\}$ select maximal subtrees avoiding $r$, with coefficient
$-1/2$. They partition the intersection-two neighbors, each with its unique
shared pair. Their proper variable sets leave a spare point beyond $r$.

\paragraph{Retained roots and doubled totals.}
Retain every full pair root and the global disjoint-tree root:
\[
 A_{ij}=\sum_{k\notin\{i,j\}}x_{\{i,j,k\}},\qquad C_* = \sum_T x_T.
\]
For each $i<h-1$, use a balanced tree on its $h-1$ incident pair roots to
retain
\[
 D_i=\sum_{j\ne i}A_{ij}=2\sum_{T\ni i}x_T.
\]
The factor two counts the two other points of each triple. These point-tree
supports overlap; they are scalar sums with their actual integer
multiplicities. Scatter the $h$ retained totals by
\[
 (Rz)_S=\begin{cases}
 \tfrac14\sum_{i\in S}D_i-\tfrac12C_*,&h-1\notin S,\\
 C_*-\tfrac14\sum_{i<h-1,\ i\notin S}D_i,&h-1\in S.
 \end{cases}
\]
On source contributions, the omitted point total equals
$3C_*-\tfrac12\sum_{i<h-1}D_i$.
Substituting gives the central coefficient $(|S\cap T|-1)/2$.
No such relation is assumed for arbitrary initial scratch.

\paragraph{Binary labels and orthonormal residuals.}
For a nondegenerate binary subspace $U$, its characteristic vector is $P_Uc$,
where $c=(1,\ldots,1)$, since $x\cdot x=x\cdot c$.
If $U\subset V$ are nondegenerate, the orthogonal difference
$V\cap U^\perp$ is nonalternating precisely when the characteristic vector of
$V$ is outside $U$. Every nondegenerate nonalternating binary space admits
an orthonormal basis: while more than one dimension remains, choose a unit
vector different from the characteristic vector. Its complement remains
nonalternating; the one-dimensional case terminates the induction.

Label each leaf by its triple line. A disjoint-tree addition with source
union $K$ has the coordinate label $F_K$. A triple-to-coordinate residual
has characteristic $1_K+t_T\ne0$; coordinate-to-coordinate residuals have
coordinate unit bases. At a selected output $F_K\subset t_S^\perp$; the
spare point gives a unit in the output residual. Selecting the whole
complement of $S$ would leave an alternating two-plane and is forbidden.
Full coordinate nodes are retained for the global total; their retirement
complement is zero and needs no unit witness.

A fixed-pair node with variable set $K$ has the orthonormal basis
$\{1_P+e_k:k\in K\}$. Its child residuals are the newly added basis vectors.
Its selected output residual has a unit at the spare point outside $P,K,r$.
Put $w_i=c+e_i$. Because $h$ is even, $w_i\cdot w_j=[i=j]$.
The full pair-root label is
$U_{ij}=\operatorname{span}\{w_k:k\ne i,j\}$: the sum of its original
star basis vectors is $c+e_i+e_j$, converting each basis vector into $w_k$.
Thus its complement has basis $w_i,w_j$.
Label each addition forming $D_i$ by
\[
 S_i=w_i^\perp=\operatorname{span}\{w_k:k\ne i\}.
\]
Two distinct incident pair-root labels span $S_i$, with unit-line residuals;
later joins keep this label. Its characteristic is $e_i$ and its complement
is the unit line $w_i$. These witnesses cover activation, retirement and all
new internal residuals, including nodes whose source union is full.

\paragraph{Exact reversible scalar wrapper.}
Use the outgoing-use/pivot compiler on the complete DAG, sharing identical
source inputs and allocating distinct physical fanout roles. With $C$
additions and $q$ designated outputs it has $R_{\rm aux}=C+q$ roles.
Let $L$ be its invertible mixer, $V$ the source copy, and $J$ the side
injection. The coefficient identities give $(R+J)LV=I$.
The chronological schedule
\[
 L,\ -R,\ -J,\ L^{-1},\ V,\ L,\ R,\ J,\ L^{-1},\ -V
\]
subtracts $(R+J)Lz$, adds $(R+J)L(z+Vx)$, changes the target by $x$, and
restores every arbitrary initial scratch value. Invert the whole schedule
for the inverse shear. Forward, inverse with banks exchanged, forward gives
$(X,Y)\mapsto(-Y,X)$ on every data pair, fixing every auxiliary role.

\paragraph{Common scatter frames in both orientations.}
At stage $j$ write the original prefix decomposition $A=B\perp P$, where
$P$ and the future factor $Q$ are norm-one tensor lines. Put
\[
 D_0=(B\otimes F)\otimes Q,\quad D_1=(A\otimes F)\otimes Q,\quad
 D_U=D_0\perp(P\otimes U\otimes Q).
\]
Early mixer and scatter gates use $D_0$; late cleanup uses $D_1$.
The forward source copy uses source-line frames, the middle mixer uses
$D_U$, the entire retained-total scatter uses the common frame $D_0$, and
only then the side injection uses each target-complement frame.
Consequently each $D_i$ output descends once from $D_{S_i}$ to $D_0$,
losing $h-1$ dimensions; the global output descends from $D_1$ to $D_0$,
losing $h$. The scatter is an actual scalar gate touching all totals and
data targets at one common frame; no decrease is free.
The reverse scatter uses $D_1$, followed by the complemented inverse mixer
at $D_{U^\perp}$, giving the same orthogonal differences. The only local
loss in either orientation is
\[
 \ell=(h-1)^2+h=553.
\]
The scalar target updates $R,J$ commute, but their distinct frame assignments
are essential. There is no direct edge $S_i\to t_S^\perp$.
Tensoring local residual bases by the unit generators of $P,Q$ preserves
orthonormality. The additional data-boundary residuals are
$B\otimes t_S^\perp\otimes Q$; coordinate units outside the earlier tensor
support and the target triple give norm-one witnesses whenever nonzero.
The original data-stage boundaries are unchanged.

\paragraph{Shared banks, endpoints and phase contract.}
Pair adjacent ground points and let $\pi(T)$ flip each point to its partner.
This triple involution has intersection zero or two with its input.
Match every complete stage-1 auxiliary bank $(A,B)$ to the stage-3 bank
$(B,\pi(A))$, keeping local physical role indices. Scalar restoration makes
this valid for arbitrary shared scratch. The joining labels are
\[
 E=F\otimes U_A\otimes U_B,\qquad
 H=(U_B\otimes U_{\pi(A)})^\perp\otimes F.
\]
Here $U_T$ is the triple line. The middle pairing vanishes, hence $E\subset H$.
A tensor of coordinate units outside $B$ in the first and third factors,
and any coordinate unit in the second, lies in $H\cap E^\perp$ and has norm
one. Its residual has dimension $m-2h$. Replacing the two separate endpoint
extensions of dimension $m-h$ by this join removes one role and exactly $m$
kernel calls. Every surviving auxiliary still has source zero and sink full.

Data terminals remain $X:U_a\to\mathcal F$ and $Y:0\to U_a^\perp$,
with $u_a$ of weight $27$. Orthogonal Hamming weight is additive modulo four,
so each edge phase difference is one forward or inverse $aI+bX_v$ kernel
per orthonormal residual basis vector. The original source/sink $Z$ corrections,
$i^{27f}$ scalar, output sign removal and bank exchange give
$C^{\otimes mf}$ on every role for every positive number $f$ of columns.
This is the pinned manuscript's binary phase argument applied to the residuals
above, with no additional recursive child calls.

\paragraph{Exact counts and recurrence.}
The disjoint and intersection-two output counts are $q_0=205956$, $q_2=27600$.
The rectangle count independently gives $123855$ additions and $205956$
output uses. Each fixed-pair tree on $n=h-2$ variables has
$q_P=n(n-1)-\sum_{\rm internal}(n-|K|)$ outputs.
Retention activates all $v-1$ global-tree additions, all
$\binom h2(h-3)$ pair additions and $(h-1)(h-2)$ point additions, or $8325$
additions before rectangles. Thus $C=132180$,
$q=q_0+q_2+h=233580$, and $R_{\rm aux}=365760$.
The $h$ totals are included in this count; no separate center bank is added.
\[
 W=2N+2v^2R_{\rm aux}=3013310215168,\qquad
 L_{\rm loss}=3v^2\ell=6796219584,
\]
\[
 s=Wm-2N+2L_{\rm loss}=41655997423981952,\qquad
 \eta=\frac{Wm-s}{Wm}=\frac{365}{5084246016}.
\]
\begin{proposition}\label{prop:retained-tree-complex-interface}
The retained-total complex motif has the all-role scalar and binary phase
contracts above, with exactly $s$ directional-kernel calls and recurrence
saving $1-\sigma=750/10^{11}$.
\end{proposition}
Indeed exact enclosure gives $\log m<477/50$ and
\[
 \eta-\frac{750}{10^{11}}\frac{477}{50}
 =\frac{11935523}{49650840000000000}>0;
\]
$e^{-x}>1-x$ therefore gives $s/W<m^{1-750/10^{11}}$.
There are $6696869422848<12W$ scalar gates, with coefficient magnitude at most
one and denominator at most four. Saved-input evaluation costs at most two
halvings and one addition per term, hence $3W^2$ operations per gate.
Including inverse-child corrections,
\[
 36W^3+4s+4W+4<64(W+m+1)^3=E,
\]
and $2\le s<m^5$. The retained generalized stopped-depth guard applies with
$B=s+E$ and $C_1=5-4\beta+\zeta$.

\paragraph{Compact-control and assembly transfer.}
Instantiate the existing compact row reservations with this $m,W,s$.
The totals are scalar roles; they introduce no address coordinates.
Every recursive child still has exactly $V/W$ volume and all complete
compact fields. The written movement, repair, reservation and recurrence
proofs give the same internal, leaf and reservation exponents
\[
 \chi=\tau+(1-\beta)\max\{\sigma-\tau,0\},\quad
 \chi_{\rm leaf}=\sigma+\beta(1-\sigma),\quad
 \chi_{\rm reserve}=\max\{1-c,0\}.
\]
Keep $\tau=1-296/10^{11}$, $\epsilon=1999/10000$, $\beta=1/1000$,
$\zeta=1/10000$, $\delta=10^{-6}$ and $C_1=49961/10000$, and choose
\[
 c=1,\quad\lambda=1-2959/10^{12},\quad
 \lambda'=1-2958/10^{12},\quad\kappa=591/10^{12}>2^{-31}.
\]
All completed-layer inequalities and seven assembly margins are strict, with
\[
 \min_i g_i=g_3=\frac{2956521}{5\cdot10^{15}},\qquad
 \min_i g_i-\kappa=\frac{1521}{5\cdot10^{15}}>0.
\]
Thus the retained assembly conditionally gives
$T(n)=O(n(\log n)^{1-591/10^{12}})$, with saving $591/83$ times the preceding
$83/10^{12}$ witness. The fixed finite alphabet, fixed number of tapes,
compact-control tape proofs and pinned analytic/multiplication interfaces
remain premises. Exact finite checks do not prove the complete machine theorem.
